\documentclass[11pt]{article}
\usepackage[margin=1in]{geometry}
\usepackage{amsmath,amssymb,amsthm}
\usepackage{array}
\usepackage{parskip}

\newtheorem{theorem}{Theorem}
\newtheorem{lemma}[theorem]{Lemma}
\newtheorem{proposition}[theorem]{Proposition}
\theoremstyle{definition}
\newtheorem{definition}[theorem]{Definition}
\theoremstyle{remark}
\newtheorem{remark}[theorem]{Remark}

\title{Disproof of conjecture \texttt{00000001186}}
\author{earthking11}
\date{2026-09-12}

\begin{document}
\maketitle

\section*{Verdict: FALSE}

We disprove conjecture \texttt{00000001186} as stated. The disproof consists of three
independent internal contradictions. It is unconditional: it does not rely on the
classification of finite simple groups, nor on any structural theory of finite
groups.

\section{The conjecture}

\begin{definition}[as stated]
Let $m(k)$ denote the minimal order of a nonabelian simple group with exactly $k$
distinct prime factors.
\end{definition}

\begin{quote}
\textbf{Conjecture.} $m(3) = 60$ ($\mathrm{A}_5$), $m(4) = 504$
($\mathrm{PSL}_2(8)$), $m(5) = 660$ ($\mathrm{PSL}_2(11)$); and
$m(k)/m(k-1) \le 4$ for all $k$, with equality at $k = 4$.
\end{quote}

Throughout, $\omega(n)$ denotes the number of \emph{distinct} primes dividing $n$,
and $\Omega(n)$ denotes the number of primes dividing $n$ counted \emph{with
multiplicity}. The definition of ``exactly $k$ distinct prime factors'' for a group
$G$ means $\omega(|G|) = k$: a group with exactly $k$ distinct prime factors is one
whose order is divisible by exactly $k$ distinct primes.

\section{The three contradictions}

\subsection*{Contradiction 1: $504$ has three distinct prime factors, not four}

\begin{lemma}
$504 = 2^{3}\cdot 3^{2}\cdot 7$, hence $\omega(504) = 3$.
\end{lemma}

\begin{proof}
Direct computation: $2^{3} = 8$, $3^{2} = 9$, and $8 \cdot 9 \cdot 7 = 72 \cdot 7 =
504$. The prime divisors are $2$, $3$, and $7$, which are pairwise distinct, so
$\omega(504) = 3$.
\end{proof}

\begin{theorem}
No group of order $504$ has exactly four distinct prime factors. Consequently, on the
definition above, $m(4) \ne 504$.
\end{theorem}

\begin{proof}
Let $G$ be a group with $|G| = 504$. The distinct prime factors of $G$ (as the
definition uses the phrase) are the distinct primes dividing $|G| = 504$. By
Lemma 1 these are exactly $2, 3, 7$, so $\omega(|G|) = 3$. In particular
$\omega(|G|) \ne 4$: no group of order $504$ has exactly four distinct prime
factors. Since $m(4)$ is by definition the order of such a group, the value $504$
cannot equal $m(4)$.
\end{proof}

\begin{remark}
This contradiction is unconditional. It does not require knowing whether a
nonabelian simple group of order $504$ exists, nor which one is intended: for
\emph{every} group of order $504$, the order has only three distinct prime divisors,
whereas the definition of $m(4)$ demands exactly four. Clause (1) above stands on
its own.
\end{remark}

\subsection*{Contradiction 2: $660$ has four distinct prime factors, not five}

\begin{lemma}
$660 = 2^{2}\cdot 3 \cdot 5 \cdot 11$, hence $\omega(660) = 4$.
\end{lemma}

\begin{proof}
$2^{2} = 4$ and $4 \cdot 3 \cdot 5 \cdot 11 = 12 \cdot 55 = 660$. The prime divisors
$2, 3, 5, 11$ are pairwise distinct, so $\omega(660) = 4$.
\end{proof}

\begin{theorem}
No group of order $660$ has exactly five distinct prime factors. Consequently, on the
definition above, $m(5) \ne 660$.
\end{theorem}

\begin{proof}
If $|G| = 660$, then the distinct prime factors of $G$ are the distinct primes
dividing $660$, namely $2, 3, 5, 11$; thus $\omega(|G|) = 4 \ne 5$. Hence no group of
order $660$ has exactly five distinct prime factors, and $m(5) \ne 660$. Equivalently,
$660$ is a candidate value for $m(4)$, not for $m(5)$: the statement has mislabelled
it by one index.
\end{proof}

\subsection*{Contradiction 3: the ratio bound fails on the conjecture's own numbers}

\begin{theorem}
Using the values listed in the conjecture,
\[
\frac{m(4)}{m(3)} = \frac{504}{60} = \frac{42}{5} = 8.4 > 4 .
\]
This contradicts the asserted inequality $m(k)/m(k-1) \le 4$ for all $k$, and in
particular contradicts the assertion that equality holds at $k = 4$.
\end{theorem}

\begin{proof}
Reduce the fraction: $504/60 = (504/12)/(60/12) = 42/5$. Since
$42/5 = 8.4$ and $8.4 > 4$, the ratio exceeds $4$. The conjecture asserts
$m(k)/m(k-1) \le 4$ for \emph{all} $k$; taking $k = 4$ and substituting the stated
values $m(4) = 504$, $m(3) = 60$ gives $8.4 \le 4$, a false numerical statement.
\end{proof}

\begin{remark}
Contradiction 3 is pure arithmetic on the numbers supplied by the conjecture itself.
It is independent of what those numbers are taken to name and of the meaning assigned
to ``prime factors''.
\end{remark}

\section{Factorisation table}

\begin{center}
\begin{tabular}{r|l|c|c}
$n$ & factorisation & $\omega(n)$ (distinct) & $\Omega(n)$ (multiplicity) \\
\hline
$60$  & $2^{2}\cdot 3\cdot 5$          & $3$ & $4$ \\
$504$ & $2^{3}\cdot 3^{2}\cdot 7$      & $3$ & $6$ \\
$660$ & $2^{2}\cdot 3\cdot 5\cdot 11$  & $4$ & $5$ \\
\end{tabular}
\end{center}

Thus $\omega(504) = 3 \ne 4$ (Contradiction 1), $\omega(660) = 4 \ne 5$
(Contradiction 2), and $504/60 = 8.4 > 4$ (Contradiction 3). Each of the three
contradictions alone falsifies the conjecture; they are independent in the sense that
removing any one still leaves the conjecture false by the other two.

\section{The multiplicity reading rescues nothing}

A reader might attempt to read ``exactly $k$ distinct prime factors'' as
``$\Omega(|G|) = k$'', i.e.\ counting prime factors with multiplicity. This does not
save the conjecture. From the table:

\begin{itemize}
\item $\Omega(60) = 2 + 1 + 1 = 4 \ne 3$. Under the multiplicity reading the very
first listed value, $m(3) = 60$, already fails.
\item $\Omega(504) = 3 + 2 + 1 = 6 \ne 4$, so $m(4) = 504$ still fails.
\item $\Omega(660) = 2 + 1 + 1 + 1 = 5$. This is the only listed number that becomes
consistent under the multiplicity reading, but it is assigned to $m(5)$; it would make
$660$ a candidate for $m(5)$, not for $m(4)$.
\end{itemize}

Moreover the ratio bound is a statement about the listed numbers themselves:
$504/60 = 8.4 > 4$ holds under either reading. Therefore no reading of ``prime
factors'' makes all of the listed values simultaneously correct, and the conjecture is
false as stated.

\section{True values on the standard reading (context only)}

The disproof above does not need the following; we record it as context. On the
standard reading, using $\omega$, the initial values are
\[
m(3) = 60 \;(\mathrm{A}_5), \qquad m(4) = 660 \;(\mathrm{PSL}_2(11)).
\]
Indeed $660 = 2^{2}\cdot 3\cdot 5\cdot 11$ has exactly four distinct prime divisors,
and every nonabelian simple group of order strictly less than $660$ has at most three
distinct prime divisors. The orders occurring below $660$ are
\[
60 = 2^{2}\cdot 3\cdot 5,\quad
168 = 2^{3}\cdot 3\cdot 7,\quad
360 = 2^{3}\cdot 3^{2}\cdot 5,\quad
504 = 2^{3}\cdot 3^{2}\cdot 7,
\]
each with exactly three distinct prime divisors; here
$|\mathrm{A}_5| = 60$, $|\mathrm{PSL}_2(7)| = 168$, $|\mathrm{A}_6| = 360$, while
$504$ is not the order of any nonabelian simple group. Hence, on the standard reading,
the statement's ``$m(4) = 504$'' is wrong, and its ``$m(5) = 660$'' is a one-index
mislabelling of the correct $m(4)$.

We stress that this section is context only. The refutation rests solely on the three
contradictions, of which Contradictions 1 and 2 require only the factorisations
$504 = 2^{3}\cdot3^{2}\cdot7$ and $660 = 2^{2}\cdot3\cdot5\cdot11$, and Contradiction 3
requires only $504/60 = 8.4 > 4$. In particular the refutation does not depend on
classifying simple groups.

\end{document}
